%% IntentDiff-RAG — Elsevier CAS double-column
\documentclass[a4paper,fleqn]{cas-dc}
\usepackage[authoryear,longnamesfirst]{natbib}

\begin{document}
\let\WriteBookmarks\relax
\def\floatpagepagefraction{1}
\def\textpagefraction{.001}

\shorttitle{}
\shortauthors{}

\title[mode=title]{}

\author[1]{}
\cormark[1]
\ead{}
\credit{}

\affiliation[1]{organization={},
            addressline={},
            city={},
            postcode={},
            country={}}

\cortext[1]{Corresponding author}

\begin{abstract}
hallo
\end{abstract}

\begin{highlights}
\item
\item
\item
\end{highlights}

\begin{keywords}
\sep
\end{keywords}

\maketitle

\section{Pendahuluan}\label{sec:intro}

% Setiap klaim faktual di bagian ini sudah dicocokkan ke sumber aslinya
% (lihat docs/references.md). Jangan menambah angka tanpa sumber terverifikasi.

Pengembangan perangkat lunak modern bergantung pada paket pihak ketiga yang
didistribusikan melalui repositori terpusat. Di ekosistem JavaScript, repositori
npm menerima sekitar 15 pembaruan paket setiap menit, dan pemasangan satu paket
dapat menarik puluhan dependensi transitif \citep{wyss2025evaluating}. Setiap
versi baru dari paket-paket tersebut sampai ke proyek hilir melalui mekanisme
yang sama: versi diterbitkan, lalu dipasang oleh proyek yang bergantung padanya.
Akibatnya, kepercayaan terhadap sebuah paket pada praktiknya juga berarti
kepercayaan terhadap setiap versi yang kelak diterbitkan atas nama paket itu.

Kepercayaan inilah yang dimanfaatkan penyerang. Alih-alih menerbitkan paket baru
yang harus lebih dulu ditemukan korbannya, penyerang menyisipkan kode berbahaya
ke dalam pembaruan paket yang sudah digunakan secara luas, yaitu serangan yang
oleh \citet{sofaer2024rogueone} disebut \emph{rogue update}. Rangkaian insiden
terbaru menunjukkan besarnya risiko ini. Pada September 2025, penyerang
memperoleh kredensial seorang pengelola (\emph{maintainer}) melalui rekayasa
sosial, lalu memperbarui 18 paket populer, termasuk \texttt{chalk} dan
\texttt{debug}. Muatan yang disisipkan secara diam-diam memanipulasi transaksi
mata uang kripto, sementara paket-paket tersebut secara gabungan diunduh sekitar
dua miliar kali per minggu \citep{aws2025chalk}. Sepekan kemudian, cacing
\emph{Shai-Hulud} mulai mencuri token npm dan kredensial lain dari lingkungan
pengembang. Dengan token curian itu, ia menerbitkan versi terinfeksi dari
paket-paket lain milik pengelola yang sama, dan lebih dari 180 paket
terkompromi dalam 24 jam pertama \citep{aws2025chalk,unit42shaihulud}. Pada Maret
2026, akun salah satu pengelola \texttt{axios} dibajak dan dipakai untuk
menerbitkan versi 1.14.1 dan 0.30.4. Kedua rilis itu menambahkan dependensi
\texttt{plain-crypto-js}, yang saat instalasi mengunduh dan menjalankan
\emph{remote access trojan} \citep{datadog2026axios}. Dalam ketiga kasus, muatan
berbahaya tiba melalui kanal distribusi yang sah milik paket yang sudah
dipercaya.

Mekanisme pengamanan yang lazim dipakai pengembang tidak dirancang untuk
menghadapi serangan semacam ini. Perintah \texttt{npm audit}, misalnya,
mengirimkan deskripsi dependensi proyek ke registry dan meminta laporan
kerentanan yang sudah diketahui \citep{npmaudit}, padahal pembaruan berbahaya
yang baru diterbitkan belum memiliki catatan semacam itu. Deteksi karenanya
harus dilakukan terhadap isi pembaruan itu sendiri, sebelum ada laporan publik.
Kesulitannya, kode berbahaya dan kode yang sah menggunakan API yang sama.
Berdasarkan evaluasi berskala besar terhadap detektor paket berbahaya npm,
\citet{guo2026effective} menyimpulkan bahwa setiap detektor harus menyelesaikan
ambiguitas antara kemampuan (\emph{capability}) kode dan niat (\emph{intent}) di
baliknya. Sebagai contoh, alat bantu \emph{build} yang sah seperti
\texttt{husky}, \texttt{patch-package}, dan \texttt{prisma} mendeklarasikan
\emph{install hook} yang sama dengan yang dipakai malware. Ambiguitas ini mahal
secara operasional. Karena pembaruan yang sah jauh lebih banyak daripada yang
berbahaya, tingkat positif palsu yang kecil pun dapat menghasilkan alarm palsu
yang jumlahnya melebihi deteksi yang benar \citep{wyss2025evaluating}.

Penelitian sebelumnya menempuh dua pendekatan. Pendekatan pertama adalah analisis
program. RogueOne membandingkan aliran data antara versi lama dan versi baru
sebuah paket, lalu menandai aliran data baru yang melintasi \emph{trust domain}.
Pada dataset pembaruan tunggal, pendekatan ini mendeteksi 75,6\% \emph{rogue
update} dengan 4,7\% positif palsu \citep{sofaer2024rogueone}. Namun, analisis
berbasis interpretasi abstrak ini membutuhkan rata-rata 226,7 detik per
pembaruan, belum termasuk kasus yang melampaui batas waktu satu jam. Konfigurasi
yang lebih sensitif memang menaikkan deteksi menjadi 92,7\%, tetapi dengan 40\%
positif palsu \citep{sofaer2024rogueone}. Pendekatan kedua memanfaatkan
\emph{large language model} (LLM). SocketAI memakai GPT-3 dan GPT-4 untuk
meninjau kode paket npm; dengan GPT-4, sistem ini mencapai presisi 99\% dan F1
97\% \citep{zahan2025leveraging}. \citet{wyss2025evaluating} memusatkan analisis
LLM pada pembaruan paket. Dengan \emph{smart-diff} yang mereduksi ukuran masukan
sebesar 95,4\%, GPT-4o mendeteksi seluruh 46 pembaruan berbahaya historis dan
163 klon berbahaya, dengan 8 positif palsu dari 2.000 pembaruan tipikal.
IntelGuard menambahkan \emph{retrieval-augmented generation} (RAG) dengan basis
pengetahuan yang dibangun dari 8.024 halaman laporan intelijen ancaman, lalu
menilai apakah semantik kode sesuai dengan fungsi yang dimaksudkan paket. Pada
paket npm, sistem ini mencapai akurasi 98,72\% dengan 0,20\% positif palsu
\citep{guo2026bridging}.

Meskipun menjanjikan, hasil-hasil tersebut menyisakan tiga keterbatasan.
\emph{Pertama}, keberhasilan deteksi berbasis LLM sejauh ini ditunjukkan dengan
model tertutup: GPT-4o, GPT-4-Turbo, dan Gemini 2.0 Flash pada
\citet{wyss2025evaluating}, serta GPT-4.1-mini pada IntelGuard
\citep{guo2026bridging}. Model tertutup mengharuskan kode dikirim ke penyedia
model, dan penggunaan LLM secara ekstensif menimbulkan biaya finansial maupun
energi yang tidak kecil \citep{wyss2025evaluating}. Evaluasi model
\emph{open-weight} yang dapat dijalankan secara lokal memang sudah ada, tetapi
dilakukan pada tingkat paket utuh dengan model berukuran 6,7 miliar parameter
\citep{nguyen2025taint}. Kemampuan model berukuran 3--4 miliar parameter dalam
menilai \emph{smart-diff} pembaruan belum pernah diukur.

\emph{Kedua}, \emph{retrieval} pada IntelGuard hanya mengambil contoh berbahaya
\citep{guo2026bridging}. Padahal, literatur \emph{in-context learning}
menunjukkan bahwa pemilihan dan urutan contoh memengaruhi prediksi model. Model
bahasa cenderung memprediksi label yang sering muncul atau yang muncul di dekat
akhir \emph{prompt} \citep{zhao2021calibrate}, dan contoh yang dipilih
berdasarkan kedekatan dengan masukan menghasilkan kinerja yang lebih baik
daripada contoh acak \citep{liu2022makes}. Karena itu, konteks yang hanya berisi
contoh berbahaya berpotensi mendorong model menjatuhkan vonis berbahaya, termasuk
pada pembaruan sah yang perilakunya mirip. Contoh positif dan negatif telah
diusulkan untuk memperjelas karakteristik keluaran yang diinginkan dan yang
harus dihindari model \citep{gao2024customizing}. Namun, sejauh penelusuran
kami, pendekatan ini belum dievaluasi untuk deteksi paket berbahaya. Selain itu,
temuan bahwa mengganti label contoh secara acak hampir tidak menurunkan kinerja
\citep{min2022rethinking} menunjukkan bahwa contoh pembanding perlu membawa alasan
yang eksplisit, bukan sekadar label.

\emph{Ketiga}, gagasan membandingkan perilaku program dengan deskripsi yang
dinyatakan pengembangnya bukanlah hal baru. CHABADA mengelompokkan aplikasi
Android berdasarkan deskripsinya, lalu menandai penggunaan API yang menyimpang
dari kelompoknya. Tanpa memerlukan pola malware yang sudah diketahui,
pendekatan ini menandai 56\% malware baru \citep{gorla2014checking}. IntelGuard
menerapkan penalaran serupa secara implisit \citep{guo2026bridging}. Efek
profil intent yang eksplisit --- deskripsi, README, dan \emph{scripts} dari
versi sebelumnya --- terhadap positif palsu pada pembaruan sah yang
menambahkan perilaku sensitif belum diukur melalui ablasi yang terkontrol.

Penelitian ini menempuh ketiga keterbatasan itu melalui IntentDiff-RAG,
sebuah \emph{pipeline} analisis statis untuk menilai pembaruan paket npm dengan
LLM \emph{open-weight} berukuran kecil. Untuk setiap pasangan versi,
IntentDiff-RAG melakukan empat langkah:
(i)~membentuk \emph{smart-diff} mengikuti spesifikasi \citet{wyss2025evaluating};
(ii)~menyusun profil intent dari versi sebelumnya, karena isi versi baru berada
di bawah kendali penyerang \citep{wyss2025evaluating};
(iii)~mengambil satu contoh berbahaya dan satu contoh sah yang mirip
berdasarkan kesamaan kemampuan yang ditambahkan oleh pembaruan; dan
(iv)~meminta model menghasilkan vonis dalam format JSON beserta baris kode yang
menjadi bukti. Pipeline ini dievaluasi pada dua model, yaitu
Qwen3-4B-Instruct-2507 dan Qwen2.5-Coder-3B-Instruct. Himpunan datanya
dipisahkan menjadi tiga subset: pembaruan berbahaya, pembaruan tipikal, dan
\emph{hard negative}, yaitu pembaruan sah yang menambahkan kemampuan sensitif
dengan tujuan yang terdokumentasi. Setiap perbandingan konfirmatori hanya
berbeda pada satu faktor, sehingga efek contoh pembanding dan efek profil intent
dapat dipisahkan dari efek jumlah contoh.

Kontribusi penelitian ini adalah sebagai berikut.

\begin{itemize}
\item Evaluasi LLM \emph{open-weight} berukuran 3--4 miliar parameter untuk
mendeteksi pembaruan npm berbahaya pada tingkat \emph{smart-diff}, dengan metrik
yang dilaporkan terpisah untuk pembaruan tipikal dan \emph{hard negative}.
\item Desain ablasi yang memisahkan efek contoh sah pembanding dari efek jumlah
contoh, serta efek profil intent dari efek \emph{retrieval}. Aturan keputusan
dan koreksi uji berganda ditetapkan sebelum eksperimen dijalankan.
\item Protokol data yang dapat direplikasi, meliputi subset \emph{hard negative}
berlabel ganda, skema \emph{knowledge card}, seleksi sampel yang deterministik,
dan pengelompokan kampanye untuk mencegah kebocoran antara basis pengetahuan dan
data uji.
\item Pengukuran biaya inferensi lokal, yaitu token, latensi median dan persentil
ke-95, serta kebutuhan memori GPU, untuk setiap komponen \emph{pipeline}.
\end{itemize}

% TODO(hasil): tambahkan satu paragraf ringkasan temuan utama setelah eksperimen
% selesai. Jangan diisi sebelum angka hasil tersedia.

Bagian selanjutnya disusun sebagai berikut. Bagian~\ref{sec:related} membahas
penelitian terkait. Bagian~\ref{sec:method} menjelaskan IntentDiff-RAG.
Bagian~\ref{sec:setup} memaparkan data, konfigurasi, dan prosedur evaluasi.
Bagian~\ref{sec:results} menyajikan hasil, Bagian~\ref{sec:discussion} membahas
implikasi dan ancaman validitas, dan Bagian~\ref{sec:conclusion} menutup artikel.

\section{Penelitian Terkait}\label{sec:related}

% Semua angka di bagian ini dicocokkan ke sumber aslinya (docs/references.md).

Bagian ini membahas enam kelompok penelitian yang menjadi dasar IntentDiff-RAG:
serangan rantai pasok pada ekosistem paket, deteksi paket berbahaya tanpa LLM,
deteksi pembaruan berbahaya, deteksi berbasis LLM, penggunaan \emph{retrieval}
dan contoh dalam \emph{prompt} untuk analisis keamanan, serta pemeriksaan
kesesuaian antara deskripsi dan perilaku program. Bagian ini ditutup dengan
posisi penelitian ini terhadap keenam kelompok tersebut.

\subsection{Serangan rantai pasok pada ekosistem paket}

Risiko rantai pasok di npm berakar pada struktur ekosistemnya.
\citet{zimmermann2019small} menunjukkan bahwa sebuah paket npm secara implisit
mempercayai rata-rata 79 paket pihak ketiga dan 39 pengelola. Mereka juga
menemukan bahwa segelintir akun pengelola cukup untuk menyuntikkan kode
berbahaya ke sebagian besar paket. \citet{ohm2020backstabber} mengumpulkan dan
menganalisis 174 paket berbahaya yang dipakai dalam serangan nyata di npm, PyPI,
dan RubyGems antara November 2015 dan November 2019. \citet{ladisa2023sok}
kemudian menyusun taksonomi serangan rantai pasok perangkat lunak
\emph{open-source} berbentuk \emph{attack tree}. Taksonomi ini memuat 107 vektor
serangan yang dikaitkan dengan 94 insiden nyata dan 33 langkah mitigasi.

Penelitian ini hanya menyasar satu kelas serangan, yaitu pembaruan berbahaya pada
paket yang sebelumnya sah. Pembatasan yang sama dipakai RogueOne, yang
mengecualikan paket berbahaya baru seperti \emph{typosquatting} dari cakupannya
\citep{sofaer2024rogueone}. Karena banyak serangan bekerja saat instalasi,
sebagian penelitian memilih mencegah alih-alih mendeteksi. Latch, misalnya,
membuat \emph{manifest} izin yang merangkum perilaku setiap paket saat
instalasi, lalu memeriksanya terhadap kebijakan yang ditetapkan pengguna
\citep{wyss2022wolf}.

\subsection{Deteksi paket berbahaya tanpa LLM}

Detektor awal menggabungkan beberapa jenis analisis. MalOSS memadukan analisis
metadata, statis, dan dinamis. Dari 339 paket berbahaya baru yang dilaporkan ke
pengelola registry, 278 dikonfirmasi \citep{duan2021towards}. Amalfi
menggabungkan \emph{classifier} berbasis \emph{machine learning}, pemeriksa
reproduksibilitas paket, dan pendeteksi klon. Ekstraksi fitur dan klasifikasinya
hanya membutuhkan beberapa detik per paket \citep{sejfia2022practical}.

Penelitian berikutnya memodelkan perilaku sebagai urutan pemanggilan API. DONAPI
menggabungkan analisis statis dengan urutan pemanggilan API dari eksekusi
dinamis, lalu mencocokkannya dengan basis pengetahuan perilaku yang berjenjang
\citep{huang2024donapi}. Cerebro memakai satu model urutan perilaku berbahaya
untuk npm dan PyPI sekaligus, dan menemukan 196 paket berbahaya baru di npm serta
306 di PyPI \citep{zhang2025killing}. OSCAR menjalankan paket di dalam
\emph{sandbox}, melakukan \emph{fuzzing} terhadap fungsi dan kelas yang diekspor,
serta memantau perilakunya melalui \emph{hook} API. Pada npm, OSCAR mencapai F1
0,95 dan menurunkan tingkat positif palsu dari 34,63\% menjadi 2,57\%
\citep{zheng2024towards}.

Evaluasi terpadu oleh \citet{guo2026effective} membandingkan 11 alat (16 varian)
pada 6.420 paket berbahaya dan 7.288 paket sah. Dari evaluasi tersebut, mereka
menyimpulkan bahwa presisi dan \emph{recall} setiap alat ditentukan oleh cara
alat itu menyelesaikan ambiguitas antara kemampuan dan niat kode.

Keterbatasan kelompok penelitian ini bagi masalah yang kami teliti ada dua.
Pertama, sebagian besar detektor di atas menilai paket secara utuh, bukan
perubahan antarversi. Kedua, detektor dinamis memerlukan eksekusi paket, padahal
penelitian ini membatasi diri pada analisis statis.

\subsection{Deteksi pembaruan berbahaya}

Kelompok penelitian ini menjadikan perubahan antarversi sebagai unit analisis.
RogueOne menandai aliran data baru yang melintasi \emph{trust domain} antara
versi lama dan baru. Pendekatan ini mendeteksi 75,6\% \emph{rogue update} dengan
4,7\% positif palsu pada dataset pembaruan tunggal \citep{sofaer2024rogueone}.
Analisis RogueOne juga menunjukkan titik buta pendekatan diferensial. Ketika
muatan berbahaya hanya memakai API yang sudah digunakan versi sebelumnya, tidak
muncul aliran data baru, sehingga konfigurasi bawaannya melewatkan seluruh 10
kasus serangan tersamar dalam kategori tersebut \citep{sofaer2024rogueone}.

\citet{froh2023differential} menyatakan perilaku mencurigakan sebagai kueri
CodeQL dan membandingkan hasil kueri antarversi. Pendekatan ini mengenali semua
versi berbahaya pada dataset paket yang memiliki riwayat kode berbahaya, dengan
rata-rata 1,3\% alarm palsu pada paket npm populer yang sah.

\citet{wyss2025evaluating} membawa unit analisis yang sama ke LLM melalui
\emph{smart-diff}. Analisis mereka menunjukkan bahwa positif palsu terutama
berasal dari operasi sensitif, penghapusan kode dalam jumlah besar (termasuk
paket yang diganti dengan penanda \emph{deprecation}), dan perubahan dependensi.
Kombinasi obfuskasi aliran data ringan dan pemecahan muatan juga dapat
menyebabkan negatif palsu.

\subsection{LLM untuk deteksi paket berbahaya}

SocketAI menunjukkan bahwa GPT-4 mengungguli analisis statis berbasis CodeQL,
dengan peningkatan presisi sebesar 16\% dan F1 sebesar 9\%, pada 5.115 paket npm
yang 2.180 di antaranya berbahaya \citep{zahan2025leveraging}. MalPacDetector
memanfaatkan LLM untuk membangkitkan fitur secara otomatis alih-alih mengandalkan
fitur yang dirumuskan pakar. Hasilnya adalah tingkat positif palsu 1,3\% dan
negatif palsu 7,5\%, serta 39 paket berbahaya baru yang dikonfirmasi tim keamanan
npm \citep{wang2025malpacdetector}.

IntelGuard mengambil contoh berbahaya yang mirip secara semantik dari basis
pengetahuan berisi 4.420 konteks kode berbahaya, lalu menggunakan GPT-4.1-mini
untuk menilai kesesuaian perilaku kode dengan fungsi paket
\citep{guo2026bridging}. Dalam benchmark \citet{guo2026effective}, pendekatan
yang mendasarkan penilaian pada bukti hasil \emph{retrieval} ini mencapai F1
95,98\%. Sebaliknya, SocketAI yang meminta model menilai kode tanpa bukti
tambahan hanya mencapai \emph{recall} 44,41\%.

Untuk model \emph{open-weight}, \citet{nguyen2025taint} memotong kode berdasarkan
alur \emph{taint} menuju API sensitif sebelum menganalisisnya dengan
DeepSeek-Coder 6,7B. Pendekatan ini mencapai akurasi 87,04\% pada 2.537 paket.
Pada ekosistem PyPI, evaluasi terhadap 13 LLM menemukan bahwa model mengenali
pola kode yang umum, tetapi melewatkan indikator yang memerlukan konteks lebih
luas atau pemahaman tentang niat penulis \citep{ryan2026evaluation}.

\subsection{\emph{Retrieval} dan contoh dalam \emph{prompt} untuk analisis keamanan}

Penelitian deteksi kerentanan memberi dua pelajaran yang relevan bagi desain
contoh pembanding. Pelajaran pertama: membedakan kode bermasalah dari kode yang
mirip tetapi aman merupakan tugas yang sulit bagi LLM. Vul-RAG menunjukkan bahwa
LLM hanya mencapai akurasi 0,06--0,14 dalam membedakan kode rentan dari kode
hasil tambalan yang mirip. Dengan pengetahuan kerentanan multidimensi, yaitu
semantik fungsional, akar penyebab, dan cara perbaikan, akurasinya meningkat
16--24\% \citep{du2026vulrag}. PrimeVul memperlihatkan bahwa benchmark yang kurang
realistis melebih-lebihkan kinerja: model 7B yang mencapai F1 68,26\% pada BigVul
hanya mencapai 3,09\% pada PrimeVul \citep{ding2025vulnerability}.

Pelajaran kedua: manfaat \emph{retrieval} bagi model \emph{open-weight} tidak
dapat diasumsikan. Replikasi Vul-RAG dengan beragam model \emph{open-weight}
menemukan bahwa akurasi berpasangan tertahan di sekitar 0,30, dan penambahan
kapasitas model saja tidak banyak memperbaikinya \citep{kaniewski2026revisiting}.

Di luar domain keamanan, literatur \emph{in-context learning} menunjukkan bahwa
pemilihan, urutan, dan label contoh memengaruhi prediksi
\citep{liu2022makes,zhao2021calibrate,min2022rethinking}. Selain itu, contoh
positif dan negatif telah diusulkan untuk memperjelas keluaran yang diinginkan
\citep{gao2024customizing}. Temuan-temuan ini menjadi dasar keputusan kami untuk
mengontrol jumlah dan urutan contoh, alih-alih mengasumsikan bahwa
\emph{retrieval} pasti membantu.

\subsection{Kesesuaian deskripsi dan perilaku}

Gagasan memeriksa apakah perilaku program sesuai dengan deskripsinya berkembang
di ekosistem aplikasi seluler. WHYPER menggunakan pemrosesan bahasa alami untuk
menemukan kalimat dalam deskripsi aplikasi yang menjelaskan alasan aplikasi
membutuhkan suatu izin \citep{pandita2013whyper}. AutoCog mengukur kesesuaian
antara deskripsi dan izin. Pada sebelas izin, sistem ini mencapai presisi
rata-rata 92,6\% dan \emph{recall} 92,0\%, lalu diterapkan pada 45.811 aplikasi
\citep{qu2014autocog}. CHABADA mengelompokkan lebih dari 22.500 aplikasi Android
berdasarkan topik deskripsinya, lalu menandai penggunaan API yang menyimpang di
setiap kelompok \citep{gorla2014checking}. Pada paket npm, IntelGuard menerapkan
penalaran serupa secara implisit \citep{guo2026bridging}. Namun, belum ada
penelitian yang mengukur efek profil intent eksplisit secara terpisah dari
komponen lain.

\subsection{Posisi penelitian ini}

Tabel~\ref{tab:positioning} merangkum posisi IntentDiff-RAG terhadap penelitian
yang paling dekat. Penelitian ini mengambil unit analisis pembaruan dari RogueOne,
\citet{froh2023differential}, dan \citet{wyss2025evaluating}, serta mengambil
gagasan \emph{retrieval} dan penalaran berbasis fungsi paket dari IntelGuard.
Perbedaannya terletak pada tiga aspek yang diuji secara terkontrol:
\begin{enumerate}
  \item penggunaan model \emph{open-weight} berukuran 3--4 miliar parameter pada
  tingkat pembaruan;
  \item contoh sah pembanding sebagai pasangan dari contoh berbahaya; dan
  \item profil intent eksplisit yang diambil dari versi sebelumnya.
\end{enumerate}
Mengikuti pelajaran dari Vul-RAG, PrimeVul, dan replikasinya, evaluasi dilakukan
pada \emph{hard negative}, yaitu pembaruan sah yang menambahkan kemampuan
sensitif, bukan hanya pada pembaruan tipikal.

\begin{table*}[t]
\caption{Posisi IntentDiff-RAG terhadap penelitian terdekat. Tanda ``--''
berarti komponen tersebut tidak digunakan atau tidak dijelaskan dalam sumber
yang kami periksa. Kolom \emph{Retrieval} merujuk pada pengambilan contoh ke
dalam \emph{prompt} LLM.}
\label{tab:positioning}
\small
\begin{tabular}{p{3.2cm} p{1.6cm} p{3.9cm} p{3.2cm} p{2.4cm} p{1.8cm}}
\hline
Penelitian & Unit analisis & Pendekatan & Model bahasa & \emph{Retrieval} & Intent \\
\hline
MalOSS \citep{duan2021towards} & Paket & Metadata, statis, dinamis & -- & -- & -- \\
Amalfi \citep{sejfia2022practical} & Paket & \emph{Classifier}, reproduksi, deteksi klon & -- & -- & -- \\
DONAPI \citep{huang2024donapi} & Paket & Statis dan urutan API dinamis & -- & -- & -- \\
OSCAR \citep{zheng2024towards} & Paket & Dinamis (\emph{sandbox}, \emph{fuzzing}) & -- & -- & -- \\
\citet{froh2023differential} & Pembaruan & Statis diferensial (CodeQL) & -- & -- & -- \\
RogueOne \citep{sofaer2024rogueone} & Pembaruan & Aliran data diferensial & -- & -- & -- \\
SocketAI \citep{zahan2025leveraging} & Paket & LLM & Tertutup (GPT-3, GPT-4) & -- & -- \\
\citet{wyss2025evaluating} & Pembaruan & \emph{Smart-diff} + LLM & Tertutup (GPT-4o, GPT-4-Turbo, Gemini 2.0 Flash) & -- & -- \\
IntelGuard \citep{guo2026bridging} & Paket & RAG + LLM & Tertutup (GPT-4.1-mini) & Contoh berbahaya & Implisit \\
\citet{nguyen2025taint} & Paket & \emph{Taint slicing} + LLM & \emph{Open-weight} (6,7B) & -- & -- \\
\textbf{IntentDiff-RAG} & \textbf{Pembaruan} & \textbf{\emph{Smart-diff} + RAG + LLM} & \textbf{\emph{Open-weight} (3--4B)} & \textbf{Berbahaya + sah pembanding} & \textbf{Eksplisit} \\
\hline
\end{tabular}
\end{table*}

\section{Method}\label{sec:method}

This section formalises the detection task, then describes the research
workflow (Fig.~\ref{fig:workflow}) and the system architecture that implements
it (Fig.~\ref{fig:architecture}). Not every configuration uses every
component: without examples, the input is the \emph{smart-diff} alone or the
\emph{smart-diff} plus the intent profile; in the ablations, retrieved examples
are added to the prompt. The data, models, configurations, and metrics are
specified in Section~\ref{sec:setup}.

\subsection{Problem formulation}\label{sec:problem}

The unit of analysis is a single update of a package, called a \emph{pair}:
\begin{equation}
p = \bigl(\texttt{pkg}@\mathrm{v}_{t-1},\; \texttt{pkg}@\mathrm{v}_{t}\bigr).
\label{eq:pair}
\end{equation}
Let $V(\texttt{pkg})$ be the published versions of the package, $\prec$ the
semver order, and $\tau(\mathrm{v})$ the publication time of version
$\mathrm{v}$. The previous version is the highest semver version that is both
lower than $\mathrm{v}_t$ and published before it:
\begin{equation}
\mathrm{v}_{t-1} = \max_{\prec}\,\bigl\{\mathrm{v} \in V(\texttt{pkg}) :
\mathrm{v} \prec \mathrm{v}_t \;\wedge\; \tau(\mathrm{v}) < \tau(\mathrm{v}_t)\bigr\}.
\label{eq:prev}
\end{equation}
Publication time alone is not used: a package with several release lines can
publish versions on different lines almost simultaneously, such as
\texttt{axios@1.14.1} and \texttt{axios@0.30.4} in the same incident
\citep{datadog2026axios}, and a time-only rule would compare two unrelated
lines. As in RogueOne and \citet{wyss2025evaluating}, the object under
assessment is a change to a previously legitimate package, not a package that
was malicious from its first release.

\paragraph{Capabilities.}
Let $\mathcal{S}$ be a fixed vocabulary of ten static signals
(Table~\ref{tab:signals}): eight \emph{capabilities} detected on a single
version, $C(\mathrm{v}) \subseteq \mathcal{S}_{\mathrm{cap}}$, and two
\emph{change signals} computed from both versions,
$G(\mathrm{v}_{t-1}, \mathrm{v}_t) \subseteq \mathcal{S}_{\mathrm{chg}}$.
The capability delta of a pair is
\begin{equation}
\Delta(p) = \bigl(C(\mathrm{v}_t) \setminus C(\mathrm{v}_{t-1})\bigr)
\;\cup\; G(\mathrm{v}_{t-1}, \mathrm{v}_t).
\label{eq:delta}
\end{equation}
$\Delta(p)$ is the retrieval key only; it is never shown to the model.

\paragraph{Model input.}
For a configuration $c$, the model receives a prompt
\begin{equation}
x_c(p) = \Pi\bigl(\tilde{D}(p),\; \mathrm{ver}(p),\; \mathbb{1}_c^{\mathrm{int}} \cdot I(\mathrm{v}_{t-1}),\; R_c(p)\bigr),
\label{eq:prompt}
\end{equation}
where $\tilde{D}(p)$ is the truncated \emph{smart-diff}
(Section~\ref{sec:smartdiff}), $\mathrm{ver}(p)$ the version change (e.g.\
\texttt{2.3.0}$\,\rightarrow\,$\texttt{2.4.0}), $I(\mathrm{v}_{t-1})$ the intent
profile (Section~\ref{sec:intent}) included only when
$\mathbb{1}_c^{\mathrm{int}} = 1$, $R_c(p)$ the ordered list of zero to two
retrieved cards (Section~\ref{sec:retrieval}), and $\Pi$ the fixed prompt
template. The package name, the label, the source, and the campaign are not
part of $x_c(p)$.

\paragraph{Verdict.}
The model $f_\theta$ returns a text $o = f_\theta(x_c(p))$. A parser
$\phi$ maps it to $\{\texttt{malicious}, \texttt{benign}, \bot\}$, where $\bot$
denotes an invalid output (Section~\ref{sec:verdict}). The binary prediction is
\begin{equation}
\hat{y}_c(p) =
\begin{cases}
1, & \phi(o) = \texttt{malicious},\\
0, & \phi(o) = \texttt{benign},\\
\beta, & \phi(o) = \bot,
\end{cases}
\label{eq:verdict}
\end{equation}
with $\beta = 0$ in the main analysis and $\beta = 1$ in the sensitivity
analysis. $\hat{y} = 1$ means the update is classified as malicious.

\subsection{Research workflow}\label{sec:workflow}

\begin{figure*}[t]
\centering
\IfFileExists{figures/workflow.pdf}{%
  \includegraphics[width=\textwidth]{figures/workflow.pdf}%
}{%
  \fbox{\parbox[c][6cm][c]{0.95\textwidth}{\centering
    Placeholder: \texttt{figures/workflow.pdf}\\(export from \texttt{figures/workflow.drawio})}}%
}
\caption{Research workflow. Gate G1 fixes which comparisons are tested
inferentially from the labelled counts; Gate G2 allows a single revision of
the output format on the dev set and then freezes the design before any run on
the test set.}
\label{fig:workflow}
\end{figure*}

The study proceeds in four phases separated by two gates
(Fig.~\ref{fig:workflow}).

\emph{Data and labels.} Malicious updates (MAL), typical updates (TYP),
\emph{hard negatives} (HN), and the knowledge-base pool are collected
(Section~\ref{sec:data}). Two annotators label every MAL and HN candidate
independently, and Cohen's~$\kappa$ is computed. At Gate~G1, the labelled count
of each subset is compared with its minimum to decide which comparisons are
tested inferentially and which are reported descriptively; the decision is
recorded before any model is run.

\emph{Preprocessing.} Every pair is unpacked and turned into a
\emph{smart-diff}, a capability delta, and an intent profile inside a container
without network access. Annotators write the knowledge-base cards, and the
retriever selects cards for each pair and configuration.

\emph{Design freeze.} All seven configurations are run on a ten-pair dev set to
check output validity, truncation, and retrieval exclusion. At Gate~G2, one
revision of the output format is allowed if the outputs are not valid
(Section~\ref{sec:verdict}); after that, the prompts, configurations, signal
vocabulary, and decision rules are frozen. No prompt is tuned against test-set
results.

\emph{Evaluation.} Stage~1 runs the seven configurations on both models at
temperature~0. Stage~2 repeats two configurations with three sampling seeds.
Metrics, statistical tests, and an error analysis follow
(Section~\ref{sec:metrics}).

\subsection{System architecture}\label{sec:architecture}

\begin{figure*}[t]
\centering
\IfFileExists{figures/architecture.pdf}{%
  \includegraphics[width=\textwidth]{figures/architecture.pdf}%
}{%
  \fbox{\parbox[c][7cm][c]{0.95\textwidth}{\centering
    Placeholder: \texttt{figures/architecture.pdf}\\(export from \texttt{figures/architecture.drawio})}}%
}
\caption{System architecture. The host only fetches and checksums tarballs;
unpacking and all analysis of package content run in a container without
network access. The capability delta reaches the model only through
retrieval. Annotation tooling produces the gold labels and the
knowledge-base cards, and annotators never see model output.}
\label{fig:architecture}
\end{figure*}

The system has four parts (Fig.~\ref{fig:architecture}): acquisition, isolated
analysis, inference, and annotation. Table~\ref{tab:components} lists the
components. They exchange files rather than calls: each stage reads the
previous stage's output, so every intermediate artefact can be inspected and
every stage re-run independently.

\emph{Acquisition} runs on the host. The fetcher queries the npm registry and
OSV advisories, downloads tarballs, verifies each against the registry's
\texttt{dist.integrity} (SHA-512) or \texttt{dist.shasum}, and records its
SHA-256. Malicious versions come from the Datadog dataset, the Wyss et al.\
replication package, or RogueOne artefacts.

\emph{Isolated analysis} runs in a container started with
\texttt{--network none}. Tools such as UglifyJS and Webcrack read sample code
as data; network isolation guarantees that no request reaches an endpoint
found in a sample, even if a tool behaves unexpectedly. Every entry point
refuses to run if any network interface other than loopback is present.
Malicious samples are analysed statically only: they are never installed,
executed, imported, or opened as URLs.

\emph{Inference} sends each prompt to a local vLLM server one request at a
time, so that latency is not distorted by overlapping requests, and stores the
raw output with the parsed verdict, token counts, latency, and the identity,
label, score, and position of every retrieved card.

\emph{Annotation} is a web application used for label audit, card writing, and
HN verification. Each annotator's decisions are stored in a separate file and
the interface never shows another annotator's labels or any model output, so
the labels stay independent and are not influenced by predictions.

\begin{table}[t]
\caption{Components of the system and their responsibilities.}
\label{tab:components}
\small
\begin{tabular}{p{2.3cm} p{5.3cm}}
\hline
Component & Responsibility \\
\hline
Fetcher & Downloads tarballs on the host, checks registry integrity, records SHA-256. \\
Unpacker & Extracts both versions of a pair inside the isolated container. \\
\emph{Smart-diff} generator & Normalises JavaScript, diffs the two trees, truncates to the token budget. \\
Capability extractor & Detects the signals of Table~\ref{tab:signals} on raw files and computes $\Delta(p)$. \\
Intent profiler & Builds $I(\mathrm{v}_{t-1})$ from \texttt{package.json} and the README of $\mathrm{v}_{t-1}$. \\
Knowledge base & Approved cards written by annotators; never model output. \\
Retriever & Ranks cards by Jaccard similarity of deltas and selects $R_c(p)$. \\
Prompt builder & Assembles $x_c(p)$ from the fixed template (Eq.~\ref{eq:prompt}). \\
Inference server & vLLM with \texttt{bfloat16} weights, sequential requests, temperature~0. \\
Parser & Maps the raw output to a verdict or $\bot$ (Eq.~\ref{eq:verdict}). \\
Evaluator & Computes metrics, tests, tables, and figures from the stored results. \\
Annotation tool & Audit, card writing, and HN verification; one decision file per annotator. \\
\hline
\end{tabular}
\end{table}

\subsubsection{\emph{Smart-diff}}\label{sec:smartdiff}

The \emph{smart-diff} $D(p)$ follows the generator in the replication package
of \citet{wyss2025evaluating}, which reduced input size by 95.4\% in the
original study. Files whose names end in \texttt{js} are minified
(\texttt{uglifyjs -c -m}) and then unminified with Webcrack in both versions,
so that variable names and comments do not become spurious signals. The two
trees are then compared with \texttt{diff -N -U 0 -r} after removing
\texttt{*.wasm}, \texttt{*.md}, \texttt{*.map}, \texttt{.npmignore}, licence
files, \texttt{*.tsbuildinfo}, and lock files. \texttt{package.json} is
included; the README is not. The original pipeline also removes repeated code
with Dolos ($k{=}70$, window 17), but the script it calls is not part of the
replication package, so this step is currently omitted.

Two additions cover what the source does not specify. First, lines longer
than 2{,}000 characters are cut to 2{,}000 characters with a marker of the
removed length: a single encoded payload line consumes the budget without
adding information after its beginning, and its obfuscation signal is already
captured by the extractor on the raw files. Second, the diff is truncated once
per pair to $B = 28{,}000$ tokens with the Qwen2.5-Coder-3B-Instruct
tokenizer,
\begin{equation}
\tilde{D}(p) = \mathrm{trunc}_B\bigl(D(p)\bigr),
\label{eq:trunc}
\end{equation}
and the same truncated diff is used by both models and every configuration.
The budget follows from the 32{,}768-token context (Table~\ref{tab:models})
minus the output, system prompt, intent profile, two cards, and template
margin. The mean \emph{smart-diff} in \citet{wyss2025evaluating} is 25{,}246
tokens and 1.47\% exceed 128{,}000 tokens, so a truncation rule is required.
$\mathrm{trunc}_B$ splits the diff per file and adds files in priority
order---the \texttt{package.json} diff, then files that carry extractor
signals in $\mathrm{v}_t$ (most signals first), then the remaining files by
path---cutting the last file at a hunk boundary. A truncated pair carries a
closing marker and the flag \texttt{truncated}, and is analysed as a subgroup.

\subsubsection{Capability extractor}\label{sec:extractor}

The extractor reads the \emph{raw} files of both versions, before
minification, because the \emph{smart-diff} transformation renames
identifiers and can erase obfuscation patterns. It scans JavaScript files and
\texttt{package.json} with text patterns (Table~\ref{tab:signals}). The
vocabulary describes capabilities rather than malware rules: contrastive
retrieval must find legitimate updates with the same capability, and detector
rules are designed to match malicious patterns only. For the eight
capabilities, the delta holds signals detected in $\mathrm{v}_t$ but not in
$\mathrm{v}_{t-1}$ at package level; \texttt{install\_script} also enters the
delta when an existing hook's command changes. For a card without a previous
version, the delta is every detected capability.

\begin{table}[t]
\caption{Signal vocabulary $\mathcal{S}$. The last two are change signals
computed from both versions.}
\label{tab:signals}
\small
\begin{tabular}{p{2.1cm} p{5.5cm}}
\hline
Signal & Detected when \\
\hline
\texttt{network} & \texttt{http}, \texttt{https}, \texttt{net}, \texttt{dgram}, \texttt{tls}, or \texttt{dns} is imported, or \texttt{fetch}, \texttt{XMLHttpRequest}, \texttt{WebSocket}, or \texttt{sendBeacon} is called \\
\texttt{process\_exec} & \texttt{child\_process} is imported \\
\texttt{host\_info} & \texttt{process.env} is read whole, or OS identity is read (\texttt{userInfo}, \texttt{hostname}, \texttt{homedir}, \texttt{networkInterfaces}) \\
\texttt{fs\_sensitive} & a string literal names a credential path such as \texttt{.ssh}, \texttt{.npmrc}, \texttt{.aws}, or \texttt{/etc/passwd} \\
\texttt{browser\_data} & cookies, forms, web storage, password fields, or submit listeners are accessed \\
\texttt{dynamic\_code} & \texttt{eval}, \texttt{new Function}, or the \texttt{vm} module runs code \\
\texttt{obfuscation} & a long base64 or hex literal, \texttt{\_0x} identifiers, \texttt{Buffer.from} with base64 or hex, or \texttt{atob} \\
\texttt{install\_script} & \texttt{preinstall}, \texttt{install}, or \texttt{postinstall} in \texttt{scripts} \\
\texttt{dependency\_change} & a new dependency, or a dependency moved to a non-registry source \\
\texttt{large\_deletion} & at least 50\% of the text files of $\mathrm{v}_{t-1}$ are gone in $\mathrm{v}_t$ \\
\hline
\end{tabular}
\end{table}

\texttt{browser\_data} is included because 33 of the 46 original incidents in
\citet{wyss2025evaluating} are form grabbing; without it, most MAL pairs would
have no retrieval key. \texttt{host\_info} ignores named variables such as
\texttt{NODE\_ENV}, which are routine configuration, and targets collection of
the whole environment or the host identity. New dependencies and large
deletions are attack paths and false-positive sources in both
\citet{wyss2025evaluating} and RogueOne \citep{sofaer2024rogueone}. The 50\%
threshold for \texttt{large\_deletion} was fixed before looking at the data.
The delta is blind to payloads that reuse a capability already present in
$\mathrm{v}_{t-1}$; recall on MAL pairs with an empty delta is reported
separately (Section~\ref{sec:metrics}).

\subsubsection{Intent profile}\label{sec:intent}

The intent profile $I(\mathrm{v}_{t-1})$ is taken from the previous version
only: the threat model assumes that an attacker can change any file in an
update \citep{wyss2025evaluating}, including a README that states a false
purpose. It contains the \texttt{description} and \texttt{keywords} of
\texttt{package.json} (up to 128 tokens), every \texttt{scripts} entry (up to
256 tokens), and the beginning of the root README with images, badges, and
HTML tags removed, up to a total of 1{,}024 tokens. A missing field is written
as \texttt{not available}; a pair with neither description nor README keeps
its place in the test set and is counted. The fixed limit keeps the
\emph{smart-diff} budget identical across configurations, so the intent
ablation does not change the diff the model sees.

The package name and the capability delta are not shown. The prompt of
\citet{wyss2025evaluating} also contains only the \emph{smart-diff}; hiding the
name reduces the chance that the model recognises an incident from memorised
training data rather than from the code. A name that appears inside the diff
or the README is left in place, and its effect is estimated by separating
incidents published before and after each model's release date.

\subsubsection{Knowledge base and retrieval}\label{sec:retrieval}

Each card carries a label, the package's purpose in one sentence, a behaviour
chain (e.g.\ \emph{install hook} $\rightarrow$ \emph{read whole environment}
$\rightarrow$ \emph{HTTPS POST to a hard-coded host}), a code excerpt of at
most 15 verbatim lines from the raw files, and a rationale of at most two
sentences covering the trigger, the data and its destination, and the
relation to the stated purpose. Its package name, campaign, source, and delta
are not shown to the model. IntelGuard links code fragments to behaviour
descriptions and expert reasoning \citep{guo2026bridging}; the card fields
follow that structure. Behaviour is written as a chain because co-occurring
capabilities, such as collection followed by exfiltration, strengthen the
detection signal \citep{guo2026effective}. The rationale is explicit because
randomly replacing demonstration labels barely degrades in-context learning
\citep{min2022rethinking}, so a card carrying only a label would not suffice.
Cards are written in English from package content by annotator~A and checked
by annotator~B against the source code; model output is never a source of a
card. A card is at most 512 tokens.

The pool holds 50 malicious and 25 legitimate cards. Malicious cards come from
packages that were malicious from the start, which lie outside the update
threat model \citep{sofaer2024rogueone} and are therefore valid as knowledge
rather than as test items. Campaigns that also appear in the MAL test set are
skipped, with at most two cards per campaign. Legitimate cards are the 25 HN
pairs held out in Section~\ref{sec:data}.

For a pair $p$, the candidate set excludes cards that share the package name
or the campaign:
\begin{equation}
\mathcal{K}_p = \{\,k \in \mathcal{K} : \mathrm{pkg}(k) \neq \mathrm{pkg}(p)
\;\wedge\; \mathrm{camp}(k) \neq \mathrm{camp}(p)\,\}.
\label{eq:candidates}
\end{equation}
Each candidate is scored with the Jaccard similarity of the deltas,
\begin{equation}
J(p, k) =
\begin{cases}
\dfrac{|\Delta(p) \cap \Delta(k)|}{|\Delta(p) \cup \Delta(k)|}, & \Delta(p) \cup \Delta(k) \neq \emptyset,\\[1ex]
0, & \text{otherwise}.
\end{cases}
\label{eq:jaccard}
\end{equation}
Candidates are ranked by descending $J$; ties are broken by ascending SHA-256
of \texttt{pair\_id|card\_id}. With only ten signals, ties are frequent, and
breaking them by knowledge-base order would select the same card for every
common delta. Text embeddings are not used: ``similar'' here means similar
capabilities, which embeddings of code or README text do not guarantee, and
IntelGuard's \texttt{text-embedding-3-large} is a closed model outside the
scope \citep{guo2026bridging}. Nearest examples are better demonstrations than
random ones \citep{liu2022makes}.

$R_c(p)$ takes the top-ranked card regardless of label (E2), the top malicious
card (E3), the two top malicious cards (E3-k2), or the top malicious and the
top legitimate card (E4 and E4-noint). Retrieval never reads the intent
profile, so two configurations that differ only in the intent block always
receive the same cards in the same order. With two cards, the top malicious
card is placed at position
\begin{equation}
\mathrm{pos}(p) =
\begin{cases}
1, & \mathrm{crc32}(\texttt{pair\_id}) \bmod 2 = 0,\\
2, & \text{otherwise},
\end{cases}
\label{eq:position}
\end{equation}
and the other slot holds the second card: the second-best malicious card for
E3-k2 and the best legitimate card for E4. Example order and label frequency
bias predictions \citep{zhao2021calibrate}; this counterbalancing keeps E3-k2
and E4 different only in the second card, not in the position of the first.
Combining malicious and legitimate examples in one prompt follows
\citet{gao2024customizing}.

\subsubsection{Verdict and output parsing}\label{sec:verdict}

The model acts as a program analyser, following the role-play of
\citet{wyss2025evaluating}. Security-sensitive behaviour is not malicious by
itself: the verdict rests on the capabilities added, the execution trigger,
the data and its destination, and---when provided---the consistency with the
intent profile and the reference examples. The prompt does not name the
\emph{hard-negative} categories, so that the false-positive rate on that
subset is not made optimistic relative to the evaluation of Wyss et al.

The output is one JSON object with \texttt{evidence} (up to three lines copied
from the \emph{smart-diff}), \texttt{reason}, and \texttt{verdict}
(\texttt{malicious} or \texttt{benign}). The parser $\phi$ reads exactly one
JSON object starting at the first \texttt{\{}. The output is valid if
\texttt{verdict}, stripped and lower-cased, is \texttt{malicious} or
\texttt{benign}, \texttt{evidence} is a list of strings, and \texttt{reason}
is a string; extra keys are ignored. An output that fails these checks, or
that was cut by the token limit, is $\bot$. Following RogueOne's treatment of
errors and timeouts \citep{sofaer2024rogueone}, $\bot$ counts as
\texttt{benign} in the main analysis; the sensitivity analysis counts it as
\texttt{malicious} (Eq.~\ref{eq:verdict}), and both are reported. The raw
output is always stored.

\section{Rancangan eksperimen}\label{sec:setup}

Pengukuran dasar memakai konfigurasi tanpa contoh. Ablasi kemudian menambah
atau menukar contoh hasil retrieval, dan memisahkan efek profil intent, pada
pair, \emph{smart-diff}, dan system prompt yang sama. Yang dilaporkan adalah
\emph{recall} pada pembaruan berbahaya, FPR pada pembaruan tipikal dan pada
\emph{hard negative}, serta token, latensi, dan memori GPU. Prompt,
implementasi \emph{smart-diff}, dan sampelnya berbeda dari
\citet{wyss2025evaluating}, dan GPT-4o tidak dijalankan ulang, sehingga selisih
terhadap angka mereka tidak dapat diuji. Angka penelitian lain dibahas di
Bagian~\ref{sec:discussion} sebagai konteks, tanpa klaim lebih baik atau lebih
buruk.

\subsection{Posisi terhadap evaluasi terdahulu}

Tabel~\ref{tab:eval-protocol} membandingkan rancangan, bukan hasil. Tugas yang
diukur sama dengan \citet{wyss2025evaluating}: menilai pembaruan, bukan paket
utuh, dari \emph{smart-diff}. Nguyen et al.\ memakai model \emph{open-weight},
tetapi pada paket utuh dengan 6,7 miliar parameter dan masukan berupa potongan
\emph{taint} \citep{nguyen2025taint}. IntelGuard menilai paket utuh dengan RAG
contoh berbahaya dan model tertutup \citep{guo2026bridging}. RogueOne menilai
pembaruan tanpa LLM \citep{sofaer2024rogueone}.

Perbedaan rancangan ada tiga. Pertama, modelnya 3--4B dan dijalankan secara
lokal. Kedua, pembaruan sah tidak diperlakukan sebagai satu kumpulan.
PrimeVul menunjukkan bahwa benchmark yang didominasi kasus mudah
melebih-lebihkan kinerja: model 7B yang mencapai F1 68,26\% pada BigVul hanya
mencapai 3,09\% pada PrimeVul \citep{ding2025vulnerability}. Vul-RAG menemukan
bahwa membedakan kode bermasalah dari kode yang mirip tetapi aman adalah tugas
yang sulit bagi LLM \citep{du2026vulrag}. Karena itu FPR dilaporkan terpisah
untuk pembaruan tipikal dan \emph{hard negative}. Ketiga, retrieval tidak
hanya mengambil contoh berbahaya seperti IntelGuard
\citep{guo2026bridging}: contoh sah pembanding dan profil intent diukur
sebagai faktor terpisah.

\begin{table*}[t]
\caption{Rancangan evaluasi model 3--4B pada \emph{smart-diff}, dibanding
evaluasi terdekat. Tabel ini membandingkan protokol, bukan angka hasil.
Tanda ``--'' berarti tidak berlaku.}
\label{tab:eval-protocol}
\small
\begin{tabular}{p{2.6cm} p{1.7cm} p{2.4cm} p{2.8cm} p{3.4cm} p{2.8cm}}
\hline
Evaluasi & Unit & Masukan & Model & Pembaruan atau paket sah & Metrik primer \\
\hline
RogueOne \citep{sofaer2024rogueone}
  & Pembaruan
  & Aliran data
  & --
  & 300 typical
  & Deteksi, FPR \\
\citet{wyss2025evaluating}
  & Pembaruan
  & \emph{Smart-diff}
  & Tertutup
  & 2.000 typical
  & Sensitivitas, spesifisitas \\
SocketAI \citep{zahan2025leveraging}
  & Paket
  & Kode paket
  & Tertutup
  & Paket sah campuran
  & Presisi, F1 \\
IntelGuard \citep{guo2026bridging}
  & Paket
  & Kode + RAG
  & Tertutup
  & Paket sah campuran
  & Akurasi, FPR \\
\citet{nguyen2025taint}
  & Paket
  & \emph{Taint slice}
  & \emph{Open-weight} 6,7B
  & Paket sah campuran
  & Akurasi \\
Penelitian ini
  & Pembaruan
  & \emph{Smart-diff} + RAG
  & \emph{Open-weight} 3--4B
  & 100 typical + 100 \emph{hard negative}
  & Recall, FPR per subset \\
\hline
\end{tabular}
\end{table*}

\subsection{Data}\label{sec:data}

Semua pengacakan memakai seed 2026. Urutan deterministik adalah urutan naik
SHA-256 dari \texttt{nama@versi}. Target himpunan uji adalah sekitar 50 pair
berbahaya (minimum 40), 100 typical, dan 100 \emph{hard negative}. Setiap
kampanye diwakili maksimal tiga pair, agar satu insiden seperti Shai-Hulud
tidak mendominasi \emph{recall}.

\paragraph{Pembaruan berbahaya (MAL).}
Kandidat diambil, menurut urutan ini, dari 46 pembaruan berbahaya
\citet{wyss2025evaluating} yang berasal dari Backstabber's Knife Collection
\citep{ohm2020backstabber}, dari entri dataset Datadog yang menandai paket
terkompromi, dari laporan OSV \texttt{MAL-*} OpenSSF pada paket yang punya
versi sah sebelumnya, dan dari artefak RogueOne \citep{sofaer2024rogueone}.
Pair masuk jika $\mathrm{v}_t$ tercatat berbahaya, $\mathrm{v}_{t-1}$ ada dan
tanpa advisory, kedua tarball tersedia, dan perubahan berbahaya berada di
berkas teks. Duplikat \texttt{nama@versi} digabung. 163 klon berbahaya pada evaluasi Wyss et al.\ tidak masuk himpunan uji: di
evaluasi asli seluruh klon terdeteksi, sehingga memasukkannya akan menaikkan
\emph{recall} tanpa menambah kasus yang sebanding dengan insiden asli. Empat
klon dipakai hanya pada \emph{dev set}.

\paragraph{Pembaruan tipikal (TYP).}
Seratus pembaruan diambil acak dari 2.000 typical updates Wyss et al.: 50 dari
1.000 paket yang paling banyak dijadikan dependensi, dan 50 dari 1.000 paket
acak, pada snapshot npm 15 April 2024. Pembagian 50/50 mempertahankan campuran
paket populer dan paket biasa seperti di sumbernya. TYP diasumsikan sah,
mengikuti asumsi yang sama; pembaruan yang ternyata punya advisory diganti
dengan hasil undian berikutnya.

\paragraph{\emph{Hard negative} (HN).}
Subset ini adalah pembaruan sah yang menambahkan kemampuan sensitif dengan
tujuan yang terdokumentasi. Konstruksinya mengikuti tiga temuan yang sudah ada,
bukan kategori baru. Tabel~VI pada \citet{wyss2025evaluating} membagi positif palsu menjadi
operasi sensitif, penghapusan besar (termasuk penanda \emph{deprecation}), dan
perubahan dependensi; pada analisis langsung, contohnya meliputi validasi
lisensi yang di-obfuscate dan unggahan ke S3 di skrip uji.
\citet{guo2026effective} menunjuk \texttt{husky}, \texttt{patch-package}, dan
\texttt{prisma} sebagai sumber positif palsu karena \emph{install hook}.
RogueOne mencatat dependensi baru, skrip instalasi, dan API sistem baru sebagai
sumber alarm palsu \citep{sofaer2024rogueone}.

Lima kategori di Tabel~\ref{tab:hn} menyaring kandidat dari sisa typical
updates Wyss et al.\ dan dari pencarian registry npm. Kandidat diterima jika
tidak ada advisory, delta kemampuan tidak kosong, tujuan perilakunya
terdokumentasi, dan penerbitnya sama dengan versi sebelumnya. Kategori
ditentukan oleh tujuan yang terdokumentasi, bukan oleh sinyal penyaring.
Kuota uji adalah 20 pair per kategori (minimum 15). Dua puluh lima pair
lainnya, lima per kategori, disisihkan sebagai contoh sah di basis pengetahuan
untuk eksperimen retrieval dan tidak masuk skor FPR.

\begin{table}[t]
\caption{Kategori \emph{hard negative}. Sinyal hanya menyaring kandidat;
kategori mengikuti tujuan yang terdokumentasi.}
\label{tab:hn}
\small
\begin{tabular}{p{2.2cm} p{2.0cm} p{3.4cm}}
\hline
Kategori & Sinyal & Dasar di evaluasi terdahulu \\
\hline
Eksekusi proses dan perkakas
  & \texttt{install\_script}, \texttt{process\_exec}
  & \emph{Install hook} sah \citep{guo2026effective}; skrip instalasi pada RogueOne \\
Telemetri
  & \texttt{network}
  & Operasi sensitif, Tabel~VI Wyss et al. \\
Perizinan
  & \texttt{obfuscation}, \texttt{dynamic\_code}
  & Validasi lisensi ter-obfuscate pada analisis langsung Wyss et al. \\
Jaringan dan unduhan
  & \texttt{network}, \texttt{process\_exec}
  & Operasi sensitif; API sistem baru pada RogueOne \\
Depresiasi dan perubahan struktur
  & \texttt{large\_deletion}, \texttt{dependency\_change}
  & Penghapusan besar dan perubahan dependensi, Tabel~VI Wyss et al. \\
\hline
\end{tabular}
\end{table}

\paragraph{Kampanye dan pelabelan.}
Dua pair masuk satu kampanye jika merujuk insiden yang sama, memuat berkas
tambahan yang identik, atau kode JavaScript tambahannya mirip menurut Dolos
(parameter Wyss: $k{=}70$, jendela 17) pada ambang 0,5. Ambang condong ke
penggabungan: memisahkan item yang sebenarnya satu kampanye lebih merusak
validitas daripada menggabungkan dua kampanye yang berbeda. Dua anotator
melabeli MAL dan HN secara independen sebelum model dijalankan pada himpunan
uji. TYP tidak dilabeli ganda. Cohen's~$\kappa$ dilaporkan; ketidaksepakatan
yang tidak terselesaikan dikeluarkan dan dihitung.

\paragraph{\emph{Dev set}.}
Sepuluh pair (empat klon berbahaya Wyss et al., tiga TYP sisa, tiga HN dari
bagian yang disisihkan untuk basis pengetahuan) dipakai untuk memeriksa
parsing, pemotongan diff, dan pengecualian retrieval. Dev set tidak dipakai
untuk menyetel aturan keputusan berdasarkan akurasi.

\subsection{Model dan inferensi}\label{sec:models}

Dua model di Tabel~\ref{tab:models} mewakili satu model instruksi umum dan
satu model khusus kode pada rentang 3--4B. Keduanya
satu keluarga dengan ukuran kosakata yang sama, sehingga jumlah token
sebanding; generalisasi ke keluarga lain terbatas. Qwen2.5-Coder-3B-Instruct
memakai lisensi \texttt{qwen-research}, bukan Apache-2.0.

\begin{table}[t]
\caption{Model yang diuji. Konteks Qwen3-4B dibatasi ke 32.768 token agar
kedua model menerima masukan identik.}
\label{tab:models}
\small
\begin{tabular}{p{2.4cm} p{2.6cm} p{2.6cm}}
\hline
  & Qwen3-4B-Instruct-2507 & Qwen2.5-Coder-3B-Instruct \\
\hline
Parameter & 4,0B & 3,09B \\
Lisensi & Apache-2.0 & \texttt{qwen-research} \\
Konteks native & 262.144 & 32.768 \\
Konteks yang dipakai & 32.768 & 32.768 \\
\hline
\end{tabular}
\end{table}

Inferensi memakai vLLM, \texttt{bfloat16}, tanpa kuantisasi, permintaan
sekuensial, dan \emph{temperature} 0. Tiga permintaan pemanasan dari dev set
dibuang. Qwen3-4B secara native mendukung 262.144 token, tetapi dibatasi ke
angka Qwen2.5-Coder-3B supaya keunggulan konteks panjang tidak menjadi faktor
pembeda.

\subsection{Konfigurasi}\label{sec:configs}

Tujuh konfigurasi di Tabel~\ref{tab:configs} memakai system prompt,
\emph{smart-diff} hasil potongan, dan parameter generasi yang identik.
E0-noint hanya berisi \emph{smart-diff} dan perubahan versi: ini kondisi
terdekat dengan prompt \citet{wyss2025evaluating}. E0 menambahkan profil intent
dari $\mathrm{v}_{t-1}$. E3 mengambil satu card berbahaya teratas, mengikuti
retrieval IntelGuard yang hanya menyertakan contoh berbahaya
\citep{guo2026bridging}. E3-k2 mengambil dua card berbahaya teratas. E4
mengambil satu card berbahaya teratas dan satu card sah teratas. E4-noint
identik dengan E4 kecuali blok intent dikosongkan. E2 mengambil satu card
teratas tanpa melihat label dan tidak diuji; yang dilaporkan darinya adalah
proporsi label card yang terambil, karena pool berbahaya lebih besar (50 vs 25)
sehingga campuran KB lebih sering menghasilkan card berbahaya.

\begin{table}[t]
\caption{Konfigurasi. System prompt, \emph{smart-diff} hasil potongan, dan
parameter generasi identik di semua baris.}
\label{tab:configs}
\small
\begin{tabular}{lll}
\hline
ID & Card & Intent \\
\hline
E0-noint & -- & tidak \\
E0 & -- & ya \\
E2 & 1 teratas, label bebas & ya \\
E3 & 1 berbahaya teratas & ya \\
E3-k2 & 2 berbahaya teratas & ya \\
E4 & 1 berbahaya + 1 sah & ya \\
E4-noint & sama dengan E4 & tidak \\
\hline
\end{tabular}
\end{table}

Tiga perbandingan konfirmatori masing-masing hanya berbeda satu faktor, pada
pair dan model yang sama:

\begin{itemize}
\item E0 vs E3: ada tidaknya satu card berbahaya.
\item E3-k2 vs E4: card kedua, berbahaya terbaik kedua versus sah terbaik;
card pertama dan posisinya sama. Tanpa E3-k2, penurunan FPR pada E4 tidak
dapat dipisahkan dari efek jumlah contoh.
\item E4 vs E4-noint: ada tidaknya blok intent; card dan urutannya identik
karena retrieval tidak membaca intent.
\end{itemize}

Selisih E0 vs E0-noint, yaitu efek intent tanpa retrieval, dilaporkan
deskriptif dan tidak diuji. E2 tidak masuk perbandingan konfirmatori.

\subsection{Metrik}\label{sec:metrics}

Metrik primer dihitung per subset, bukan pada gabungan himpunan uji.

\begin{itemize}
\item \emph{Recall}: proporsi pair MAL yang diberi vonis \texttt{malicious}.
\item FPR$_{\mathrm{TYP}}$: proporsi pair TYP yang diberi vonis \texttt{malicious}.
\item FPR$_{\mathrm{HN}}$: proporsi pair HN yang diberi vonis \texttt{malicious}.
\end{itemize}

Setiap proporsi dilaporkan dengan interval kepercayaan skor Wilson 95\%, per
model dan per konfigurasi. \citet{wyss2025evaluating} memakai sensitivitas dan
spesifisitas karena himpunan berbahaya dan sah terpisah serta tidak seimbang,
sehingga presisi menjadi tidak proporsional. Alasan yang sama dipakai di sini.
Presisi dan F1 dihitung pada gabungan uji (rasio kelas 50:200 disebutkan) dan
bersifat sekunder.

FPR tidak digabung antara TYP dan HN. TYP mendekati campuran Wyss et al.\ dan
menjadi titik banding operasional --- npm menerima sekitar 15 pembaruan per
menit, sehingga setiap 1\% FPR berarti sekitar 216 alarm palsu per hari.
HN mengukur ambiguitas kemampuan versus niat yang diidentifikasi
\citet{guo2026effective}, yaitu kasus yang membuat FPR pada typical updates
terlihat rendah.

\emph{Recall} pada pair MAL dengan delta kemampuan kosong dilaporkan terpisah.
Pada pola itu payload memakai API yang sudah ada di $\mathrm{v}_{t-1}$;
konfigurasi bawaan RogueOne melewatkan seluruh 10 kasus serangan tersamar
semacam ini \citep{sofaer2024rogueone}. Subkelompok \texttt{truncated} juga
dilaporkan.

Tidak ada ambang ``cukup'' yang disepakati untuk model 3--4B. Yang dilaporkan
pada konfigurasi tanpa contoh adalah estimasi titik beserta intervalnya, per
model.

Dengan sekitar 50 pair MAL, lebar interval Wilson pada \emph{recall} 0,9
sekitar $\pm$8,5 poin; pada 100 pair HN dan FPR 0,2, lebarnya sekitar
$\pm$7,8 poin. Batas ini adalah properti desain, bukan kekurangan yang ditutup
dengan menambah sampel setelah melihat hasil.

\paragraph{Ablasi.}
Endpoint yang diuji pada tiga selisih berpasangan adalah FPR$_{\mathrm{HN}}$,
bukan F1. F1 bergantung pada rasio kelas di himpunan data; mekanisme yang
diukur --- contoh sah atau profil intent mengurangi kecenderungan menuduh ---
bekerja pada pembaruan sah yang perilakunya mirip berbahaya.

Uji memakai McNemar exact dua sisi pada subset HN, dengan luaran per pair:
diberi vonis \texttt{malicious} atau tidak. Ada enam uji (tiga selisih
$\times$ dua model), dikoreksi Holm dengan $\alpha$ keluarga 0,05. Jika tidak
ada pair yang vonisnya berbeda, $p{=}1$.

\emph{Recall} dipakai sebagai batas, bukan sebagai endpoint kedua yang diuji.
Selisih recall dihitung berpasangan. Interval 95\% berasal dari bootstrap
persentil 10.000 resample per kampanye, seed 2026, karena pair dalam satu
kampanye saling berkorelasi. Penurunan estimasi titik tidak boleh melebihi
5~poin. Dengan sekitar 50 pair MAL, batas ``tidak turun lebih dari 5~poin''
tidak dapat dibuktikan dari interval saja; yang dilaporkan adalah estimasi
titik ditambah CI.

Label hasil untuk E3-k2 vs E4 dan untuk E4 vs E4-noint:

\begin{itemize}
\item \emph{Didukung}: FPR$_{\mathrm{HN}}$ turun, $p_{\mathrm{Holm}}{<}0{,}05$,
estimasi titik penurunan recall $\le$ 5~poin, dan batas bawah CI selisih
recall $\ge{-}5$~poin.
\item \emph{Didukung sebagian}: syarat yang sama kecuali batas bawah CI
$<{-}5$~poin (recall aman belum terbukti).
\item \emph{Arah sesuai, tidak signifikan}: FPR$_{\mathrm{HN}}$ turun,
$p_{\mathrm{Holm}}{\ge}0{,}05$, dan estimasi titik penurunan recall $\le$ 5~poin.
\item \emph{Tidak didukung}: FPR$_{\mathrm{HN}}$ tidak turun, atau estimasi
titik penurunan recall $>$ 5~poin.
\end{itemize}

Untuk E0 vs E3 tidak ada syarat recall. Labelnya \emph{Didukung} jika
FPR$_{\mathrm{HN}}$(E3) $>$ FPR$_{\mathrm{HN}}$(E0) dengan
$p_{\mathrm{Holm}}{<}0{,}05$; \emph{Arah sesuai, tidak signifikan}; atau
\emph{Tidak didukung}. Klaim yang berlaku pada kedua model hanya ditulis jika
keduanya berlabel Didukung atau Didukung sebagian. ``Tidak signifikan'' bukan
bukti bahwa tidak ada efek.

\paragraph{Biaya.}
Token adalah \texttt{prompt\_tokens} dan \texttt{completion\_tokens} dari
field \texttt{usage} vLLM. Latensi adalah waktu wall-clock di sisi klien,
dari permintaan dikirim sampai respons lengkap diterima, tanpa streaming;
dilaporkan sebagai median dan persentil ke-95 (nearest-rank). Permintaan
sekuensial agar angka itu tidak terdistorsi oleh tumpang-tindih di peladen.
\citet{wyss2025evaluating} melaporkan rata-rata 9,91 detik per kueri untuk
GPT-4o.

Memori GPU adalah jumlah memori bobot yang dilaporkan log vLLM saat model
dimuat dan KV cache terhitung untuk pasangan prompt+output terpanjang yang
benar-benar terjadi. Puncak \texttt{nvidia-smi} tidak dipakai: vLLM
mengalokasikan sebagian besar VRAM saat start sesuai
\texttt{--gpu-memory-utilization}. Model GPU, total VRAM, dan dtype
dilaporkan.

Tambahan per komponen dihitung sebagai selisih: E0-noint $\rightarrow$ E0
untuk intent; E0 $\rightarrow$ E3 untuk satu card; E3 $\rightarrow$ E3-k2
atau E4 untuk card kedua. Pareto dihitung pada
(FPR$_{\mathrm{HN}}$, recall, median latensi).

\subsection{Prosedur dan apa yang tidak dibandingkan}

Urutan kerja: pelabelan selesai sebelum inferensi pada himpunan uji; dev set
dijalankan lebih dulu untuk validitas JSON; kemudian ketujuh konfigurasi di
kedua model, \emph{temperature} 0, permintaan sekuensial. Sampel berbahaya
hanya dianalisis secara statis.

FPR per kategori HN, persentase pair yang vonisnya berubah pada tiga seed
(\emph{temperature} 0,7; E0 dan E4), pemisahan insiden sebelum versus sesudah
tanggal rilis model, taksonomi kesalahan, dan kesesuaian \texttt{evidence}
dengan baris gold dilaporkan sebagai deskripsi dan tidak diuji.

Tiga perbandingan sengaja tidak dilakukan. RogueOne tidak dijalankan ulang:
rata-rata 226,7 detik per pembaruan dengan batas waktu satu jam, dan
membandingkan angka dari dataset berbeda tidak sah
\citep{sofaer2024rogueone}. GPT-4o dan GPT-4.1-mini tidak dijalankan: ruang
lingkup dibatasi model \emph{open-weight} lokal. Kalau data Wyss et al.\ di OSF
memungkinkan subset yang sama, angka kami boleh disajikan berdampingan dengan
angka mereka, tetap tanpa uji.

\section{Results}\label{sec:results}

\section{Discussion}\label{sec:discussion}

\section{Conclusion}\label{sec:conclusion}

\printcredits

\bibliographystyle{cas-model2-names}
\bibliography{cas-refs}

\end{document}
